% !TEX root = ../main.tex
\chapter[Generalized Algebraic Theories]{Generalized Algebraic Theories}

Our goal is to generate a term of a specific type in a given theory.
The goal of this chapter is to give meaning to the concept of 
term, type and theory, while introducing some tools that we will
use later when we present the algorithm.

By theories we mean Generalized Algebraic Theories (or GATs), 
introduced by John Cartmell in his thesis \cite{cartmell1978} \cite{cartmell1986},
they give us both a very idiomatic way of presenting
theories and powerful interpretation tools like
Context Categories, introduced later in this section.
The notion of generalized algebraic theory is a generalization of
the notion of many-sorted algebraic theory in the following manner:
while in Martin L{\"o}ff's many-sorted algebraic theories sorts
are constant types, interpreted as sets, in generalized
algebraic theories sorts might be constant or variable,
over a specified range of types, and variable types
are interpreted as families of sets. So a generalized
algebraic theory consists of:
\begin{enumerate}
    \item a set of sorts, that can be constant or variable;
    \item a set of operators, with the sort of the arguments and value specified;
    \item a set of axioms, where each axiom is an identity between similar expressions. 
\end{enumerate}
Let's present category theory for example:
The set of sorts contains two elements $Ob$ and $Hom$.
$Ob$ is a constant type and $Hom$ is a variable type,
for a pair of terms of type $Ob$, name them $x, y$, we
have a type $Hom(x, y)$.
Here we see a key point in general algebraic theories,
types vary over terms, no over types.
The set of operators also has two elements $id$ and $\circ$.
The $id$ operator has one argument $x$ of type $Ob$,
and returns a value of type $Hom(x, x)$. This shows
how the type of value of an operator might vary depending
on the arguments. 
The $\circ$ operator (or composition) expects 5 arguments, 
three of type $Ob$, name them $x, y, z$, one of type $Hom(x, y)$,
say $f$, and one of type $Hom(y, z)$, named $g$, and it returns
a value of type $Hom(x, z)$.
Here we see that the type of an argument might also vary depending
on a previous argument.
{\color{red}Informally, instead of writing $\circ(x, y, z, f, g)$ we can
write $\circ(f, g)$, as the first three arguments might be inferred from the last two.
This notation will be discussed in a further section.}

Now we will introduce the notation we will be using when talking 
about terms and types. {\color{orange} This notation will be formally defined
in a later section.} So consider you have a set of variables $V$,
we will associate terms to variables as required.
When asserting that a term $t$ is of type $\Delta$ we will write
$t : \Delta$, but if $t$ has variables $x_1, \dots, x_n$ occurring in it this assertion
does not make sense unless we assign some types to $x_1, \dots, x_m$.
To do this we will write:
\begin{equation*}
    \dfrac{
        x_1 : \Delta_1, \dots, x_n: \Delta_n
    }{
        t : \Delta
    }
\end{equation*}
that reads as: if $x_1$ is a variable of type $\Delta_1$, \dots and if $x_n$ is a variable of type $\Delta_n$,
then $t$ is a term of type $\Delta$. 
Similarly, if we want to assert that if $x_1$ is a variable of type $\Delta_1$,
\dots and if $x_n$ is a variable of type $\Delta_n$,
then $\Delta$ is a type, we can write:
\begin{equation*}
    \dfrac{
        x_1 : \Delta_1, \dots, x_n : \Delta_n 
    }{
        t : \Delta
    }
\end{equation*}
this assertions are called rules.

